\section{Introduction}
\label{sec:intro}

The pre-train-then-fine-tune paradigm has revolutionized modern machine learning, enabling a single high-capacity model to excel across a wide array of downstream applications through task-specific optimization \cite{dosovitskiy2020image, devlin2018bert, vaswani2017attention}. However, as the number of downstream tasks grows, maintaining, serving, and deploying separate fine-tuned expert checkpoints for each task becomes computationally prohibitive and logistically complex. To mitigate these overheads, \emph{weight-space model merging} has emerged as an attractive, zero-shot alternative for consolidating multiple fine-tuned models into a single multi-task network without requiring joint training or retraining from scratch \cite{wortsman2022model, jin2022dataless, matena2021merging}.

Despite its computational appeal, the primary obstacle to successful model merging is the phenomenon of \emph{parameter interference} and subsequent \emph{representation collapse} \cite{yadav2023ties}. When specialized expert models are fine-tuned independently from a shared pre-trained base, their parameters diverge along different, often high-dimensional and non-orthogonal trajectories. Directly interpolating or summing these task vectors (e.g., via Uniform Merging or Task Arithmetic \cite{ilharco2022editing}) often yields catastrophic performance degradation on all tasks. This failure occurs because the linear combination of unaligned, task-specific parameter updates introduces destructive interference, corrupting the specialized representation manifolds of the individual networks.

To resolve parameter interference, existing methodologies rely heavily on discrete, coordinate-wise heuristics. For instance, TIES-Merging \cite{yadav2023ties} and Sparse Task Arithmetic (STA) \cite{polymerge} employ sign voting, absolute magnitude thresholding, or hard parameter pruning (e.g., keeping only the top $50\%$ of parameters) to "denoise" the task vectors prior to combination. From a mathematical perspective, these approaches suffer from several severe limitations. First, coordinate-wise thresholding treats parameters as completely independent, ignoring the underlying low-rank geometric correlations and structural symmetries of deep neural network weights. Second, hard thresholding is non-differentiable and lacks formal mathematical guarantees or error bounds, making it impossible to reason about the representational distortion or Lipschitz properties of the resulting merged model. 

In this work, we address these fundamental limitations from the first principles of spectral theory and manifold geometry. We propose \textbf{Grassmannian Subspace Consensus Merging (GSC-Merge)}, a mathematically rigorous, partial weight-space model merging framework that targets the major linear projection layers inside the Transformer blocks (representing over 95\% of block parameters) while keeping lightweight normalization, biases, and embedding parameters task-specific to prevent statistic mismatch across highly disparate domains. Instead of heuristic coordinate-wise pruning, we formulate model merging as a problem of finding a shared consensus parameter subspace across multiple expert models. By horizontally concatenating the task vectors across all specialized experts at each layer, we construct a joint multi-task update matrix representing the collaborative update manifold. We then perform Singular Value Decomposition (SVD) on this matrix to extract the principal directions of parameter variation across all fine-tuned experts \cite{golub1996matrix}.

By projecting these joint updates onto a low-rank Grassmannian manifold $\mathbf{Gr}(r, d_{out})$, we construct a spectral filter that mathematically separates the coherent, shared consensus update directions (associated with the largest singular values) from high-frequency, task-specific orthogonal noise (associated with the tail singular values) \cite{edelman1998geometry}. Crucially, by invoking the celebrated \emph{Eckart-Young-Mirsky Theorem} \cite{eckart1936approximation}, we prove that our Grassmannian projection operator is the mathematically optimal low-rank approximation of the joint task updates under the Frobenius norm, providing a strict upper bound on the representational drift.

Furthermore, we identify and resolve a fundamental optimization challenge which we term the \emph{Overfitting-Optimizer Paradox}. When optimizing layer-wise merging coefficients on a small validation set (Offline Few-Shot Validation Tuning or OFS-Tune \cite{ofstune}), unconstrained optimization often overfits to the few-shot validation data and stream noise, resulting in poor test-set generalization. We demonstrate, both theoretically and empirically, that our Grassmannian projection acts as an elegant spectral regularizer: by restricting the parameter updates to a low-dimensional consensus manifold, it bounds the optimization space, prevents overfitting, and significantly enhances multi-task generalizability on unseen test data.

We evaluate GSC-Merge using a Vision Transformer (ViT-Tiny) backbone across four highly conflicting visual classification datasets: MNIST, FashionMNIST, CIFAR-10, and SVHN. Our empirical results validate our theoretical framework: while unconstrained validation tuning is highly sensitive to the validation split and suffers from extreme variance, GSC-Merge acts as a robust spectral regularizer that significantly reduces split-sensitivity variance (a classic bias-variance trade-off) in both the task-conditional and truly task-agnostic settings, matching the performance of unconstrained tuning while operating in a highly restricted, low-dimensional parameter space. Furthermore, GSC-Merge consistently outcompetes standard coordinate-wise heuristic baselines such as Uniform Merging, Task Arithmetic, and Sparse Task Arithmetic across all settings. To ensure scientific rigor and address common evaluation limitations, we conduct a multi-seed statistical analysis to verify the robustness of our improvements, and we provide a transparent ablation study comparing task-conditional parameter swapping with a truly task-agnostic model merging setting.

In summary, our key contributions are:
\begin{itemize}
    \item We introduce \textbf{GSC-Merge}, a mathematically principled, partial weight-space model merging framework that leverages spectral theory and Grassmannian projection to align multi-task parameter updates in major linear layers while keeping lightweight normalization and bias parameters task-specific.
    \item We provide a theoretical analysis of representation drift, proving that GSC-Merge achieves the mathematically optimal low-rank approximation of joint task updates via the Eckart-Young-Mirsky Theorem.
    \item We expose the \emph{Overfitting-Optimizer Paradox} in few-shot validation tuning and demonstrate that Grassmannian projection serves as a robust spectral regularizer to prevent optimization collapse.
    \item We conduct a rigorous 5-seed statistical analysis and evaluate GSC-Merge under both task-conditional and truly task-agnostic settings, demonstrating that our spectral consensus framework consistently outperforms coordinate-wise heuristic baselines.
\end{itemize}
